% Template file: do not edit it in a module (see "Template and module files" in README.md).
% LaTeX settings shared by the notes and the sheets (PDF output only).
% amsmath, amssymb, graphicx, xcolor and hyperref are already loaded by Quarto.
% HTML equivalents of the boxes are in style.css; keep them in sync.

% Colours (accent, alert, brand-*) are defined in generated/palette.tex, made from
% palette.yml and included before this file.

\usepackage{tcolorbox}
\tcbuselibrary{most}

% ::: {.infobox}
\newtcolorbox{infobox}{enhanced, colback=accent!4, colframe=accent!60,
  boxrule=0.5pt, arc=2pt, left=8pt, right=8pt, top=6pt, bottom=6pt,
  before skip=6pt, after skip=14pt, before upper={\setlength{\parskip}{6pt}}}

% [2]{.marks}: marks pushed to the right margin (moves to the next line if needed)
\newcommand{\marksbox}[1]{\unskip\nobreak\hfil\penalty50\hskip1em\hbox{}\nobreak
  \hfill\mbox{\bfseries\color{accent}#1}}

% [text]{.marking} and ::: {.marking}: notes for markers (staff builds only)
\newcommand{\markingfont}{\small\itshape\color{accent!85!black}}
\newtcolorbox{markingbox}{enhanced, frame hidden, sharp corners,
  colback=accent!4, borderline west={2pt}{0pt}{accent!60},
  left=6pt, right=6pt, top=3pt, bottom=3pt, before skip=4pt, after skip=6pt,
  fontupper=\markingfont, before upper={\setlength{\parskip}{4pt}}}

% [2]{.difficulty}: icons (style set by difficulty-icon in _module.yml)
\newcommand{\difficultystar}{\textcolor{accent}{$\bigstar$}}
% A module's own icon (an image file), 1em tall, slightly below the baseline
\DeclareRobustCommand{\difficultyimage}[1]{\raisebox{-0.15em}{\includegraphics[height=1em]{#1}}}
